\subsection{Springer 纤维的几何性质}

\subsubsection{连通性} Springer 纤维 $\sB_{e}$ 是连通的。参见 \cite[p.131, Prop 1]{StConj}、\cite[Ch II, Cor 1.7]{Spaltenstein} 以及习题 \ref{exI:conn}。

\subsubsection{等维性} Spaltenstein \cite{Spa77}、\cite[Ch II, Prop 1.12]{Spaltenstein} 证明：$\sB_{e}$ 的所有不可约分支维数相同。

\subsubsection{维数公式} 设 $d_{e}$ 为 $\sB_{e}$ 的维数。Steinberg \cite[Thm 4.6]{St} 与 Springer \cite{Springer} 证明
\begin{equation}\label{Spr dim}
d_{e}=\frac{1}{2}(\dim\sN-\dim G\cdot e)=\frac{1}{2}(\dim C_{G}(e)-r).
\end{equation}

\subsubsection{中心化子的作用} 群 $G\times \Gm$ 作用在 $\sN$ 上（其中 $\Gm$ 通过缩放作用）。设 $\tilG_{e}=\Stab_{G\times \Gm}(e)$ 为稳定化子，则 $\tilG_{e}$ 作用在 $\sB_{e}$ 上。注意 $\tilG_{e}$ 总是满射到 $\Gm$，其核为 $C_{G}(e)$。

$\tilG_{e}$ 在 $\cohog{*}{\sB_{e}}$ 上的作用经由有限群 $A_{e}=\pi_{0}(\tilG_{e})=\pi_{0}(G_{e})$ 分解。注意 $A_{e}$ 不仅依赖于 $G$ 的同源类，还依赖于 $G$ 的同构类。例如当 $e$ 正则时，$A_{e}=\pi_{0}(ZG)$，其中 $ZG$ 是 $G$ 的中心。$A_{e}$ 在 $\cohog{*}{\sB_{e}}$ 上的作用进一步经由 $A_{e}$ 在伴随群 $G^{\ad}$ 中的像分解。

\subsubsection{纯性} Springer \cite{SpringerPure} 证明：$\sB_{e}$ 的上同调总是纯的（当 $k=\CC$ 时按 Hodge 理论的意义，当 $k$ 为有限域时按 Frobenius 权的意义）。

\subsubsection{} 设 $e\in\sN$。考虑由包含 $\sB_{e}\incl\sB$ 诱导的限制映射 $i^{*}_{e}: \cohog{*}{\sB}\to \cohog{*}{\sB_{e}}$。则 $i^{*}_{e}$ 的像恰为 $\cohog{*}{\sB_{e}}$ 在 $A_{e}$ 下的不变量。这是 Hotta 与 Springer 的定理 \cite[Theorem 1.1]{HottaSpringer}。特别地，当 $G$ 为 $A$ 型时，这样的限制映射总是满射。

\subsubsection{奇偶消没} De Concini、Lusztig 与 Procesi \cite{DLP} 证明：对一切奇数 $i$ 与任意 $e\in\sN$，$\cohog{i}{\sB_{e}}$ 消没。当 $k=\CC$ 时他们还证明了更强的结论：奇数 $i$ 时 $\cohog{i}{\sB_{e},\ZZ}$ 消没，偶数 $i$ 时无挠。
